\documentclass[bachelor, och, autoref]{SCWorks1}
\usepackage[T1, T2A]{fontenc}
\usepackage[utf8]{inputenc}
\usepackage{graphicx}
\usepackage[sort,compress]{cite}
\usepackage{amsmath}
\usepackage{amssymb}
\usepackage{amsthm}
\usepackage{fancyvrb}
\usepackage{longtable}
\usepackage{array}
\usepackage[english,russian]{babel}
\usepackage[colorlinks=true]{hyperref}
\usepackage{float}

\begin{document}

% Кафедра (в родительном падеже)
\chair{дискретной математики и информационных технологий}

% Тема работы
\title{Разработка системы стеганографического сокрытия данных в сетевых пакетах и модуля обнаружения вредоносных вложений}

% Курс
\course{4}

% Группа
\group{421}

% Специальность/направление
\napravlenie{09.03.01 "--- Информатика и вычислительная техника}

% Фамилия, имя, отчество в родительном падеже
\author{Иванова Ивана Ивановича}

% Заведующий кафедрой
\chtitle{доцент, к.\,ф.-м.\,н.}
\chname{Л.\,Б.\,Тяпаев}

% Научный руководитель
\satitle{канд. техн. наук, доцент}
\saname{И.\,О.\,Фамилия}

% Семестр
\term{8}

% Наименование практики
\practtype{Преддипломная}

% Продолжительность практики
\duration{4}

% Даты практики
\practStart{03.02.2026}
\practFinish{11.06.2026}

% Год выполнения
\date{2026}

\maketitle


\intro

Скрытые каналы в сетевых протоколах TCP/IP активно используются APT-группировками для управления вредоносным ПО и эксфильтрации данных. Существующие средства защиты (DPI, IDS, DLP) не обеспечивают их надёжного обнаружения.

\textbf{Цель} --- разработка системы стеганографического сокрытия данных в сетевых пакетах с модулем обнаружения скрытых каналов и вредоносных вложений.

\textbf{Задачи:} анализ методов сетевой стеганографии; реализация пяти каналов (IP~ID, TCP~ISN, TCP~Timestamp, ICMP~Payload, DNS~QNAME) с криптозащитой AES-256-GCM; разработка протокола фрагментации; создание модуля обнаружения (статистика, ML, сигнатуры); анализ вредоносных вложений; тестирование.

Работа содержит 3~главы, 14~рисунков, 8~таблиц, 4~листинга. Объём --- 70~страниц.

\centerline{\textbf{КРАТКОЕ СОДЕРЖАНИЕ РАБОТЫ}}

\textbf{В первом разделе} проведён анализ ландшафта киберугроз (рост инцидентов с 1500 до 4500 за 2019--2024~гг.), применения стеганографии APT-группировками (Turla, Equation Group, APT29), нормативных требований (ФСТЭК, NIST SP~800-53). Описаны пять полей TCP/IP для стеганографии с учётом поведения ОС.

\textbf{Во втором разделе} представлены алгоритмы кодирования для каждого канала (2--1470~байт/пакет), криптозащита, протокол фрагментации NST (17-байтный заголовок, CRC32, SHA-256). Методы обнаружения: 9~статистических тестов, Isolation Forest (9~признаков), Random Forest, сигнатуры, анализ payload (20~сигнатур в 7~категориях).

\textbf{В третьем разделе} описана реализация NetStego (Python, Scapy, scikit-learn). Результаты: пропускная способность 20~Б/с -- 115{,}9~Кбит/с; F1-мера обнаружения 0{,}80--0{,}92 при FPR~$\leq$~10\,\%; восстановление данных при потере до 20\,\% пакетов. Покрытие --- 269~тестов.

\conclusion

Разработана система NetStego с пятью стеганографическими каналами TCP/IP и модулем обнаружения (статистика + ML + сигнатуры + анализ вредоносных вложений). Подтверждён компромисс скрытность/пропускная способность. Перспективы: IPv6, HTTP/2, QUIC, глубокое обучение, обнаружение в реальном времени.

\textbf{Источники:}
Wendzel S. et al., JUCS, 2016;
RFC~791, 792, 793, 1035, 7323;
NIST SP~800-53 Rev.~5.


\end{document}
